\subsection*{Position of the certified window}
\label{sec:v141-window-position}
The window parameter $\lambda$ is not the parameter used in the
published compact-window literature, which is indexed by the
logarithmic support half-width. A test supported multiplicatively in
$[\lambda^{-1},\lambda]$ is supported in $[-a,a]$ in the logarithmic
coordinate, with
\[
 a=\log\lambda=L/2 .
\]
Zhu~\cite{Zhu2026} writes this half-width $L$; it is one half of the
logarithmic length $L=2\log\lambda$ of the reading conventions. The
classical positivity range of Yoshida and Connes--Consani is stated
as $\operatorname{supp}f\subset[-(\log2)/2,(\log2)/2]$, as recalled in
\cite{Zhu2026}. On one axis:
\begin{center}
\begin{tabular}{lrr}
\toprule
result & half-width $a$ & $\lambda=e^{a}$\\
\midrule
classical Yoshida / Connes--Consani & $0.3466$ & $1.414$\\
Zhu~\cite{Zhu2026}, certified $L=0.8$ & $0.8000$ & $2.226$\\
Proposition~\ref{prop:v116-window-positive}, $\lambda=3$ & $1.0986$ & $3$\\
Proposition~\ref{prop:v126-full-window}, $\lambda=4$ & $1.3863$ & $4$\\
\bottomrule
\end{tabular}
\end{center}
The $\lambda=4$ certificate is therefore at $1.73$ times the
half-width of the published Zhu window and $4$ times the classical
range. The comparison is legitimate because the semilocal form and the
global Weil form agree on every test supported inside the window;
Proposition~\ref{prop:v121-cofinal-rh} relies on exactly this
extension-by-zero identity, so the two sides certify the same object
on the same functions.

The two results are not of the same kind. Zhu proves a certified
two-sided enclosure
$8.9\cdot10^{-18}\le\lambda_{\min}(0.8)\le2.27\cdot10^{-17}$ of the
normalized window infimum at $a=0.8$, by a one-stroke reduction to a
single finite positive semidefinite matrix.
Proposition~\ref{prop:v126-full-window} proves nonnegativity
$\mathsf W_4\succeq0$ at $a=\log4$, and
Proposition~\ref{prop:v140-ground4} encloses its complete ground,
$0<\mu_0<2.454\cdot10^{-75}$ and $\mu_1>10^{-73}$, by a structured
inverse, a simultaneous seventeen-column certificate and
direction--complement gluing. Zhu's result is stronger in kind, a
quantitative two-sided gap on a smaller window; the present one is
stronger in reach. Neither subsumes the other. The $LDL^*$ pivots and
the margin $1-\lambda_{\max}(\mathcal U,K)$ recorded in the
$\lambda=4$ certificates are not ordinary-norm eigenvalue bounds and are not comparable with
Zhu's enclosure: the pivots depend on coordinates, whereas the relative
margin compares two forms.

In the notation of Connes--Consani--Moscovici
\cite[Corollaries~3.7--3.8]{ConnesConsaniMoscovici2025}, write
$\mu^{\mathrm{CCM}}_\lambda$ for the infimum of the spectrum of the
self-adjoint operator of the form $QW_\lambda$ certified here. They
prove that $\lambda\mapsto\mu^{\mathrm{CCM}}_\lambda$ is nonincreasing,
write that they ``cannot assert that'' it is nonnegative, and prove
that $\lim_{\lambda\to\infty}\mu^{\mathrm{CCM}}_\lambda=0$ implies RH.
By that monotonicity,
\[
 \mathsf W_4\succeq0
 \quad\Longleftrightarrow\quad
 \mu^{\mathrm{CCM}}_\lambda\ge0\ \text{ for every }1<\lambda\le4,
\]
so Proposition~\ref{prop:v126-full-window} settles the sign of
$\mu^{\mathrm{CCM}}_\lambda$ on $1<\lambda\le4$, a range $1.73$ times
Zhu's certified window and beyond the windows of their own numerical
section. The floors of Proposition~\ref{prop:v138-three-floors} are the
lower bounds $\mu^{\mathrm{CCM}}_\lambda\ge-8$ for $1<\lambda\le8$.
This is a target statement, not a shortcut: cofinal nonnegativity of
$\mu^{\mathrm{CCM}}_\lambda$ is RH by
Proposition~\ref{prop:v121-cofinal-rh}, and no cofinal statement is
made here. No G2 gap or RH claim follows.
